\documentclass[10pt,a4paper]{amsart}
\usepackage[margin=19mm]{geometry}
\usepackage{amsmath,amssymb,mathtools}
\usepackage[scr=rsfs,cal=euler]{mathalfa}
\usepackage{xfrac}
\usepackage[unicode,hidelinks,bookmarks=false]{hyperref}
\newcommand{\ie}{i.e.\ }
\newcommand{\cf}{cf.\ }
\newcommand{\resp}{respectively\ }
\newcommand{\Subsection}{\subsection}
\providecommand{\nobreakdash}{\nobreak\hbox{-}}
\setlength{\emergencystretch}{2em}
\multlinegap0pt
\usepackage{bm}
\usepackage{bbm}
\usepackage{dsfont}
\usepackage{xspace}
\usepackage{cancel}
\usepackage{mathtools}
\usepackage{amsthm}
\makeatletter
\@ifundefined{thm}{\newtheorem{thm}{Thm}}{}
\@ifundefined{definition}{\newtheorem{definition}[thm]{Definition}}{}
\@ifundefined{lemma}{\newtheorem{lemma}[thm]{Lemma}}{}
\@ifundefined{notation}{\newtheorem{notation}[thm]{Notation}}{}
\makeatother
\newcommand{\C}{\mathbb{C}}
\newcommand{\N}{\mathbb{N}}
\title[Hypercube C*-algebras and an application to magic isometries]{Hypercube C*-algebras and an application to magic isometries}
\author{Bj\"orn Sch\"afer}
\date{Source: arXiv:2510.15586v1 (CC BY 4.0)}
\begin{document}
\maketitle

\noindent\emph{Source and licence.} Hypercube C*-algebras and an application to magic isometries. Version arXiv:2510.15586v1 (CC BY 4.0), distributed under the Creative Commons Attribution 4.0 International licence, as recorded in the arXiv licence metadata for this exact version and shown on the abstract page https://arxiv.org/abs/2510.15586v1. Full rights notice: licenses/arxiv-2510.15586-CC-BY-4.0.txt.\par\smallskip

\noindent\textit{Editorial scope.} Counted unit: the Lemma whose source label is \texttt{reps::lemma:decomposition\_of\_rho\_t\_for\_boundary\_points} in the article (source0.tex lines 1076--1086, source label \texttt{reps::lemma:decomposition\_of\_rho\_t\_for\_boundary\_points}), delivered together with the article's own complete original proof of it (source0.tex lines 1088--1175, one \texttt{proof} environment of 88 lines). No mathematics is written, repaired or re-derived: statement and proof are the article's own text line for line. Required context retained uncounted: pre::notation:binary\_numbers (lines 334--359); reps::def:admissible\_weightings (lines 711--729); hyp::lemma:properties\_of\_admissible\_weightings (lines 902--913); weig::def:representation\_rho\_t (lines 965--977). Every definition, display and label of those lines is retained unchanged. Omitted from this package and not redistributed: 1--333, 360--710, 730--901, 914--964, 978--1075, 1087--1087, 1176--1600. Nothing inside a retained excerpt is omitted or altered: every retained body is delivered exactly as the article's lines are written, so no omission receipt of any kind accompanies this package. Citations: the retained excerpts carry no citation command at all, so no bibliography entry of the article is transcribed and no reference is redistributed. Reference adaptation: none was needed and none was made; every reference inside the delivered unit resolves inside it, so no unresolved reference or citation is produced. Printed slips kept literal: no printed word, symbol, hypothesis or number of the counted statement or of its proof was altered. Layout and class: the article's own class is \texttt{amsart} with packages including amssymb, mathtools, hyperref, bookmark, enumitem, braket, todonotes, multicol, tikz-cd, biblatex; the wrapper is the lane's pinned \texttt{amsart} editorial wrapper at 10pt on A4, which restates the theorem-like environments the retained text needs and adds the article's own macro definitions that the retained text needs (\texttt{\textbackslash C}, \texttt{\textbackslash N}), with extra wrapper packages (bm, bbm, dsfont, xspace, cancel, mathtools). The delivered numbering is the unit's own. Assets: no retained span contains a figure environment, a graphics-inclusion command, a PDF-page inclusion, a table or a TikZ picture; no image file is copied into this package, the package declares no asset dependency and no image file is redistributed with it. Licence: version arXiv:2510.15586v1 of this article is distributed under the Creative Commons Attribution 4.0 International licence, as recorded in the arXiv licence metadata for this exact version and shown on the abstract page https://arxiv.org/abs/2510.15586v1. A full rights notice is delivered at licenses/arxiv-2510.15586-CC-BY-4.0.txt. Characters: every retained excerpt is pure ASCII, so the delivered engine, XeLaTeX with its default Latin Modern text font, reports no missing character.
\begin{notation}
    \label{pre::notation:binary_numbers}
    For every number $i \in \N$ and $n \in \N$ such that $i < 2^n$ the $n$-digit binary representation of $i$ is the sequence 
    \begin{align*}
        i_{n-1} \dots i_1 i_0 \in \{0,1\}^n \quad \text{ such that } \quad i = \sum_{k=0}^{n-1} i_k 2^{k}.
    \end{align*}
    We call the number $i_k$ the $k$-th binary digit of $i$ and denote $[i]_k := i_k$. Further, we set
    \begin{align*}
        \mathrm{par}_k(i) = \sum_{\ell=0}^{k} i_\ell \mod 2.
    \end{align*}
    Note that $\mathrm{par}_0(i)$ is the parity of the number $i$.

    Further, for every $k < n$ let $i \# k \in \N$ be the number obtained by flipping the $k$-th binary digit of $i$, i.e. 
    \begin{align*}
        i \# k = \begin{cases}
            i + 2^{k} & \text{if } i_k = 0, \\
            i - 2^{k} & \text{if } i_k = 1.
        \end{cases}
    \end{align*}
    % Finally, for an ascending sequence $\mathbf{k} = k_0 k_1 \dots k_{m-1}$ of indices with $k_{m-1} < n$ we set
    % \begin{align*}
    %     [i]_\mathbf{k} := [i]_{k_0k_1 \dots k_{m-1}} = \sum_{\ell=0}^{m} i_{k_\ell} 2^{\ell},
    % \end{align*}
    % i.e. $[i]_\mathbf{k}$ is the number obtained from $i$ by removing all binary digits except the $k_0$-th, $k_1$-th, ..., and $k_m$-th binary digit.
    % Usually, the number $n$ is clear from the context and will not be stated explicitly.
\end{notation}

\begin{definition}
    \label{reps::def:admissible_weightings}
    Let $c: E_n \to \C$ be a complex weighting of the edges of $Q_n$ and let 
    $$
    Q_n(c) = (U_n(c), V_n(c), E_n(c))
    $$
    be the subgraph of $Q_n$ induced by the edges $ij \in E_n$ with $c(ij) \neq 0$.
    We call the weighting $c$ \emph{admissible} if it satisfies
    \begin{align}
        \label{wei::eq:admissible_weighting_condition1}
        \sum_{i \in \mathcal{N}(j_1) \cap \mathcal{N}(j_2)} c(ij_1) \overline{c(ij_2)} = \delta_{j_1j_2} \qquad \text{ for all } j_1, j_2 \in V_n(c),
    \end{align}
    as well as
    \begin{align}
        \label{wei::eq:admissible_weighting_condition2}
        \sum_{j \in \mathcal{N}(i_1) \cap \mathcal{N}(i_2)} c(i_1j) \overline{c(i_2j)} = \delta_{i_1i_2} \qquad \text{ for all } i_1, i_2 \in U_n(c),
    \end{align}
    where $\mathcal{N}(x)$ with $x \in U_n \cup V_n$ denotes the set of neighbors of the vertex $x$ in $Q_n$.
\end{definition}

\begin{lemma}
    \label{hyp::lemma:properties_of_admissible_weightings}
    Let $c: E_n \to \C$ be an admissible weighting of the hypercube $Q_n$, and let $Q_n(c) \subset Q_n$ be as in Definition \ref{reps::def:admissible_weightings}. Then the following holds:
    \begin{enumerate}
        \item Assume that $Q_n(c) = \bigcup_\ell Q_n^{(\ell)}$ for disjoint sub-hypercubes $Q_n^{(\ell)} \subset Q_n$. Then the representation $\rho_c: C^\ast(Q_n) \to B(\C^{U_n(c)})$ induced by $c$ is a direct sum of irreducible representations $\rho^{(\ell)}$ which are vanishing on $p_x$ for $x \not \in Q_n^{(\ell)}$, respectively. In particular, $\rho_c$ is irreducible if and only if $Q_n(c)$ is a single sub-hypercube of $Q_n$.
        \item If $c^\prime: E_n \to \C$ is another admissible weighting of $Q_n$, then the representations $\rho_c$ and $\rho_{c^\prime}$ are unitarily equivalent if and only if $Q_n(c^\prime) = Q_n(c)$ and there are numbers $\lambda_x \in S^1$ for $x \in U_n(c) \cup V_n(c)$ such that
        \begin{align*}
            c^\prime(ij) = \lambda_i \lambda_j c(ij) \qquad \text{ for all } ij \in E_n(c).
        \end{align*}
        \item For any rank-one representation $\rho: C^\ast(Q_n) \to B(\mathcal{H})$ of the hypercube algebra $C^\ast(Q_n)$ on a Hilbert space $\mathcal{H}$ there exists an admissible weighting $c: E_n \to \C$ such that $\rho$ is unitarily equivalent to the representation $\rho_c$.
    \end{enumerate}
\end{lemma}

\begin{definition}
    \label{weig::def:representation_rho_t}
    For every tuple $\mathbf{t} = [t_0, \dots, t_{n-1}] \in \Delta_{n-1}$ let an edge weighting $c_\mathbf{t}$ of $Q_n$ be given by
    \begin{align*}
        c_\mathbf{t}(ij) = (-1)^{\mathrm{par}_k(i)} \sqrt{t_k} \qquad \text{ for all } ij \in E_n,
    \end{align*}
    where $k < n$ is the unique number with $j = i \# k$ and $\mathrm{par}_k(i)$ is the sum of the first $(k+1)$-binary digits of $i$ modulo $2$ as in Notation \ref{pre::notation:binary_numbers}. 
    Let 
    $$
        \rho_\mathbf{t}: C^\ast(Q_n) \to M_{U_n}
    $$    
    be the representation of $C^\ast(Q_n)$ induced by the weighting $c_\mathbf{t}$.
\end{definition} 

\begin{lemma}
    \label{reps::lemma:decomposition_of_rho_t_for_boundary_points}
    For $\mathbf{s} = [s_0, \dots, s_{n-2}] \in \Delta_{n-2}$ and $\ell < n$ let $\mathbf{t} \in \Delta_{n-1}$ be given by
    \begin{align*}
        \mathbf{t} = [t_0, \dots, t_{n-1}] := [s_0, \dots, s_{\ell-1}, 0, s_\ell, \dots, s_{n-2}].
    \end{align*}
    Then one has the unitary equivalence
    \begin{align*}
        \rho_\mathbf{t} \cong \left(\rho_\mathbf{s} \circ \pi^{(+\ell)}\right) \oplus \left(\rho_\mathbf{s} \circ \pi^{(-\ell)}\right).
    \end{align*}
\end{lemma}

\begin{proof}
    The representation $\rho_\mathbf{t}$ acts on $\C^{U_n}$, while the direct sum on the right-hand side of the equation in the statement acts on $\C^{U_{n-1}} \oplus \C^{U_{n-1}}$. Consider the unitary operator $\mathcal{U}: \C^{U_{n-1}} \oplus \C^{U_{n-1}} \to \C^{U_n}$ given by
    \begin{align*}
        \mathcal{U}: (e_x, 0) \mapsto e_{(x \# 0)^{(+\ell)}}, \qquad
        (0, e_x) \mapsto e_{x^{(-\ell)}},
    \end{align*}
    for all $x \in U_{n-1}$. 
    Let us determine the representation
    \begin{align*}
        \sigma(\cdot) := \mathcal{U} \left(\left(\rho_\mathbf{s} \circ \pi^{(+\ell)}(\cdot)\right) \oplus \left(\rho_\mathbf{s} \circ \pi^{(-\ell)}(\cdot)\right) \right) \mathcal{U}^\ast
    \end{align*}
    of $C^\ast(Q_n)$ on $\C^{U_n}$. For $x \in U_n$ with $[x]_\ell = 1$ we have $x=(y \# 0)^{(+\ell)}$ for a suitable (and unique) $y \in U_{n-1}$. Then, $\pi^{(+\ell)}(p_x) = \tilde{p}_y$ and $\pi^{(-\ell)}(p_x) = 0$. Thus, one easily checks
    \begin{align*}
        \sigma(p_x) = \mathcal{U} \begin{pmatrix}
            \tilde{E}_{yy} & 0 \\
            0      & 0
        \end{pmatrix} \mathcal{U}^\ast = E_{xx},
    \end{align*}
    where we write a bounded operator on $\C^{U_{n-1}} \oplus \C^{U_{n-1}}$ as a $2 \times 2$ block matrix in a natural way, and $\tilde{E}_{yy} \in M_{U_{n-1}}$ is the usual standard matrix unit. One obtains 
    \begin{align*}
        \sigma(p_x) = E_{xx}
    \end{align*}
    for $x \in U_n$ with $[x]_\ell = 0$ similarly.

    Next, let $x \in V_n$ with $[x]_\ell = 1$, and let $y \in V_{n-1}$ be the unique vertex with $x = (y \# 0)^{(+\ell)}$. Then $\pi^{(+\ell)}(p_x) = \tilde{p}_y$ and $\pi^{(-\ell)}(p_x) = 0$. Further, $\tilde{p}_y$ is the projection onto the span of the vector $\tilde{\psi}_y \in \C^{U_{n-1}}$ with entries
    \begin{align*}
        [\tilde{\psi}_y]_i = \begin{cases}
            (-1)^{\mathrm{par}_k(i)} \sqrt{s_k}     &\exists \, k < n-1 : y = i \# k, \\
            0                                       &\text{ otherwise},
        \end{cases}
    \end{align*}
    where $i$ ranges over $U_{n-1}$.
    Consequently, $\sigma(p_x)$ is the projection onto the span of the vector $\varphi_x := \mathcal{U} (\tilde{\psi}_y, 0) \in \C^{U_n}$. Since, the second component of $(\tilde{\psi}_y, 0)$ vanishes it is clear that $[\varphi_x]_i$ vanishes for $i \in U_n$ with $[i]_\ell = 0$. For the other indices $i \in U_n$ there exists a unique $\tilde{i} \in U_{n-1}$ with $i = (\tilde{i} \# 0)^{(+\ell)}$, and one computes
    \begin{align*}
        [\varphi_x]_{i} &= \begin{cases}
            (-1)^{\mathrm{par}_k(\tilde{i})} \sqrt{s_k}     &\exists \, k < n-1 : y = \tilde{i} \# k, \\
            0                                       &\text{ otherwise.}
        \end{cases}
    \end{align*}
    One checks that $y = \tilde{i} \# k$ is equivalent to
    \begin{align*}
            x &= (y \# 0)^{(+\ell)} = (\tilde{i} \# 0)^{(+\ell)} \# k = i \# k, & \text{ if } k < \ell, \\
            x &= (y \# 0)^{(+\ell)} = (\tilde{i} \# 0)^{(+\ell)} \# (k+1) = i \# (k+1), & \text{ if } k \geq \ell.
    \end{align*}
    Further,
    \begin{align*}
        \mathrm{par}_k(\tilde{i}) = \begin{cases}
            - \mathrm{par}_k((\tilde{i} \# 0)^{(+\ell)}) = -\mathrm{par}_k(i)    &\text{ if } k < \ell, \\
            \mathrm{par}_{k+1}((\tilde{i} \# 0)^{(+\ell)}) = \mathrm{par}_{k+1}(i)   &\text{ if } k \geq \ell,
        \end{cases}
    \end{align*}
    Using $\mathbf{t} = [s_0, \dots, s_{\ell-1}, 0, s_\ell, \dots, s_{n-2}]$ we thus obtain
    \begin{align*}
        [\varphi_x]_{i} &= \begin{cases}
            -(-1)^{\mathrm{par}_k(i)} \sqrt{t_k}        &\exists \, k < \ell : x = i \# k, \\
            (-1)^{\mathrm{par}_{k+1}(i)} \sqrt{t_{k+1}}       &\exists \, k \geq \ell : x = i \# (k+1), \\
            0                                           &\text{ otherwise.}
        \end{cases}
    \end{align*}
    Analogously, one checks for $x \in V_n$ with $[x]_\ell = 0$ that $\sigma(p_x)$ is the projection onto the span of the vector $\varphi_x \in \C^{U_n}$ with
    \begin{align*}
        [\varphi_x]_i = \begin{cases}
            (-1)^{\mathrm{par}_k(i)} \sqrt{t_k}                &\exists \, k < \ell : x = i \# k, \\
            (-1)^{\mathrm{par}_{k+1}(i)} \sqrt{t_{k+1}}         &\exists \, k \geq \ell : x = i \# (k+1), \\
            0                                                   &\text{ otherwise.}  \\                                
        \end{cases}
    \end{align*}
    
    Putting everything together and using $t_\ell = 0$, we obtain that $\sigma$ is induced by the edge weighting $c: E_n \to \C$ with
    \begin{align*}
        c(ij) = \begin{cases}
            \varepsilon_{j}^k (-1)^{\mathrm{par}_k(i)} \sqrt{t_k}        &\exists \, k < n : j = i \# k, \\
            0                                           &\text{ otherwise },
        \end{cases}
    \end{align*}
    for all edges $ij \in E_n$, where
    \begin{align*}
        \varepsilon_j^k = \begin{cases}
            -1,    & \text{ if } [j]_\ell = 1 \text{ and } k < \ell, \\
            1,     & \text{ otherwise.}
        \end{cases}
    \end{align*}
    Setting $\lambda_x := (-1)^{[x]_\ell \mathrm{par}_\ell(x)}$ for all $x \in U_n \cup V_n$ it is not hard to verify
    \begin{align*}
        \lambda_i \lambda_j c(ij) = c_\mathbf{t}(ij),
    \end{align*}
    where $c_\mathbf{t}$ is the edge weighting from Definition \ref{weig::def:representation_rho_t}. Thus, $\sigma$ is unitarily equivalent to $\rho_\mathbf{t}$ by Lemma \ref{hyp::lemma:properties_of_admissible_weightings}(2).
\end{proof}

\end{document}
